% Generated draft: requires mathematical and source review.
\documentclass[11pt]{article}
\usepackage[margin=25mm]{geometry}
\usepackage{fontspec}
\setmainfont{Latin Modern Roman}
\usepackage{amsmath,amssymb,mathtools,unicode-math}
\setmathfont{Latin Modern Math}
\usepackage{enumitem,needspace,xurl,hyperref}
\setlength{\parindent}{0pt}
\setlength{\parskip}{5pt}
\setlength{\emergencystretch}{3em}
\newcommand{\N}{\mathbb{N}}
\newcommand{\Z}{\mathbb{Z}}
\newcommand{\Q}{\mathbb{Q}}
\newcommand{\R}{\mathbb{R}}
\newcommand{\C}{\mathbb{C}}
\newcommand{\F}{\mathbb{F}}
\newcommand{\A}{\mathbb{A}}
\newcommand{\PP}{\mathbb{P}}
\newcommand{\E}{\mathbb{E}}
\newcommand{\eps}{\varepsilon}
\newcommand{\dd}{\,\mathrm{d}}
\begin{document}
\section*{UnsolvedMath ID: 800648}
\textbf{Source problem code:} catalogue entry
\subsection{Definitions and mathematical statement}
The supplied record does not specify the intended class of spaces or whether covers must be countable. The following gives a conditional formulation using customary open-cover definitions; these conventions require confirmation against the source.

Fix a class \(\mathcal C\) of topological spaces and a convention specifying whether only countable covers are admitted. An admissible open cover of \(X\) is a family \(\mathcal U\) of open subsets satisfying \(\bigcup\mathcal U=X\) and \(X\notin\mathcal U\), subject to that convention. Define:
\begin{itemize}
\item \(\Omega(X)\) consists of admissible open covers \(\mathcal U\) such that every finite \(F\subseteq X\) is contained in some \(U\in\mathcal U\).
\item \(\Lambda(X)\) consists of admissible open covers in which every point belongs to infinitely many members.
\item \(\Gamma(X)\) consists of infinite admissible open covers in which every point belongs to all but finitely many members.
\item \(\Omega^{gp}(X)\) consists of covers \(\mathcal U\in\Omega(X)\) admitting a partition \(\mathcal U=\bigcup_{n\geq1}\mathcal P_n\) into nonempty finite, pairwise disjoint blocks such that
\[
\forall F\subseteq X\text{ finite}\quad
\exists n_0\geq1\quad
\forall n\geq n_0\quad
\exists U\in\mathcal P_n\quad F\subseteq U.
\]
\item \(\Lambda^{gp}(X)\) consists of covers \(\mathcal U\in\Lambda(X)\) admitting such a partition with
\[
\forall x\in X\quad\exists n_0\geq1\quad
\forall n\geq n_0\quad x\in\bigcup\mathcal P_n.
\]
\end{itemize}
For classes of covers \(\mathscr A(X)\) and \(\mathscr B(X)\), the property \(\mathsf S_{\mathrm{fin}}(\mathscr A,\mathscr B)\) means that for every sequence \((\mathcal U_n)_{n\geq1}\) in \(\mathscr A(X)\), there exist finite subfamilies \(\mathcal F_n\subseteq\mathcal U_n\) such that \(\bigcup_{n\geq1}\mathcal F_n\in\mathscr B(X)\). The property \(\binom{\mathscr A}{\mathscr B}\) means that every \(\mathcal U\in\mathscr A(X)\) contains a subfamily belonging to \(\mathscr B(X)\).

\textbf{Question, conditional on the unresolved ambient conventions.} Does every \(X\in\mathcal C\) satisfy
\[
X\models\mathsf S_{\mathrm{fin}}(\Omega,\Omega^{gp})
\quad\Longleftrightarrow\quad
X\models\binom{\Omega}{\Omega^{gp}}\,?
\]
Thus the question asks whether the existence of a groupable omega subcover of every omega cover is equivalent to the stated finite-selection property.

For completeness, the supplied contextual assertion uses \(\mathsf U_{\mathrm{fin}}(\Gamma,\Gamma)\): for every sequence \((\mathcal U_n)_{n\geq1}\) of gamma covers, each having no finite subcover, there are finite \(\mathcal F_n\subseteq\mathcal U_n\) such that \(\{\bigcup\mathcal F_n:n\geq1\}\in\Gamma(X)\). The record attributes to reference label \texttt{[coc7]} the equivalence
\[
X\models\mathsf S_{\mathrm{fin}}(\Omega,\Omega^{gp})
\quad\Longleftrightarrow\quad
\forall m\geq1\;
X^m\models\mathsf U_{\mathrm{fin}}(\Gamma,\Gamma),
\]
where finite powers are understood with their product topologies. It further reports that \(\mathsf U_{\mathrm{fin}}(\Gamma,\Gamma)\) is equivalent to \(\binom{\Lambda}{\Lambda^{gp}}\). These are supplied background assertions, subject to the source's unstated hypotheses and conventions, and are not independently verified here.
\subsection{Short English statement}
Under the intended, but unspecified, conventions for spaces and open covers, ask whether every omega cover having a groupable omega subcover is equivalent to being able to select finitely many members from each cover in any sequence of omega covers so that the selected members together form a groupable omega cover. The registry labels the problem open; this draft does not independently verify that status.
\subsection{Sources}
\begin{itemize}
\item Local UnsolvedMath registry, permanent ID \(800648\), exact source problem code \texttt{catalogue entry}; title: Selection Principles Milestone Problems --- Problem 21 (\texttt{[reznicb, Problem 8]}). The supplied record lists Tsaban as proposer and \(2003\) as the proposed year, and labels the problem \texttt{open}.
\item Tsaban, \emph{Selection Principles in Mathematics: A Milestone of Open Problems} (2003), Problem 21. Supplied source URL: \url{https://arxiv.org/abs/math/0312182}. The record reports extraction by \texttt{arxiv\_latex} and access on \(2026\text{-}10\text{-}16\). This access date is supplied provenance, not an independently verified visit or literature assessment.
\item The supplied attribution/context is \texttt{[reznicb, Problem 8]}. No bibliographic expansion of this reference label is supplied.
\item The background equivalence concerning finite powers is attributed in the supplied statement to \texttt{[coc7]}. No bibliographic expansion of this reference label is supplied. The further equivalence involving \(\binom{\Lambda}{\Lambda^{gp}}\) is reported without a separate resolved citation.
\end{itemize}
\end{document}
